\documentclass[a4paper, 12pt]{article}

\usepackage{custom}


%--------------------Lettera--------------------
\def\lastversion{}
\def\title{Lettera di presentazione PB seconda parte}
\def\date{6 Giugno 2025}
%------------------------------------------------

\begin{document}

\primapagina

\newpage

\noindent
\textbf{Egregio Professor Vardanega}, \\
Con la presente, il gruppo \textit{Alt+F4} desidera candidarsi per la seconda fase della revisione di avanzamento PB (\textit{Product Baseline}) del progetto denominato \textbf{Artificial QI}, proposto dall'azienda \textbf{Zucchetti S.p.A.}\\
\\
La pagina di presentazione del gruppo, insieme alla documentazione relativa al PB, è disponibile al seguente indirizzo: 
\begin{center}
    \href{https://alt-f4-eng.github.io/Documentazione}{https://alt-f4-eng.github.io/Documentazione}\\
\end{center}
Desideriamo informarla che, nella suddetta pagina, troverà i seguenti documenti esterni del PB:

\begin{itemize}
    \item \textit{Analisi dei requisiti v.2.0}
    \item \textit{Manuale utente v.1.0}
    \item \textit{Piano di progetto v.2.0}
    \item \textit{Piano di qualifica v.2.0}
    \item \textit{Specifica tecnica v.1.0}
    \item \textit{Verbali esterni}
\end{itemize}
\noindent
E i seguenti documenti interni del PB:
\begin{itemize}
    \item \textit{Glossario v.1.0}
    \item \textit{Norme di progetto v.2.0}
    \item \textit{Verbali interni}:
    \item verbale interno del 2025/06/03
    \item verbale interno del 2025/05/26
    \item verbale interno del 2025/05/12
    \item verbale interno del 2025/05/06
    \item verbale interno del 2025/04/29
    \item verbale interno del 2025/04/22
    \item verbale interno del 2025/04/07
    \item verbale interno del 2025/03/31
    \item verbale interno del 2025/03/24
\end{itemize}
\vspace{0.75cm} 
\noindent
\textbf{Aggiornamento impegni}\\
\\
Il gruppo ha raggiunto il limite di risorse disponibili di \textbf{11.250€} prima del termine del MVP (\textit{Minimum Viable Product}), fissato per il 31/05/2025. Per questo motivo, il gruppo ha implementato unicamente la funzionalità di gestione dei dataset seguendo le prassi della buona ingegneria del software, come richiesto dal corso.\\
\vspace{0.75cm} 
\\
\noindent
Il gruppo \textit{Alt+F4} è composto dai seguenti membri:
\begin{table}[H]
    \centering
    \begin{tabular}{| l | l |}
    \hline
    \textbf{Nome Cognome} & 
    \textbf{Matricola}\\ 
        \hline
            Enrico Bianchi&
            2040978 \\
        \hline 
            Eghosa Matteo Igbinedion Osamwonyi&
            2042888 \\
        \hline 
            Guirong Lan&
            2042368 \\
        \hline 
            Pedro Leoni&
            2042359 \\
        \hline 
            Marko Peric&
            2011067 \\
        \hline 
            Francesco Savio&
            2085846 \\
        \hline 
    \end{tabular}
\end{table}
\noindent
Cordiali saluti,\\
\textbf{il gruppo \textit{Alt+F4}}

\end{document}
